\chapter{Sylow Theorems}

\section{Cauchy's Theorem and the Class Equation}

\begin{theorem}[Cauchy's theorem]\label{thm:cauchy}
Let $G$ be a finite group, and let $p$ be a prime. If
\[
    p\mid |G|,
\]
then there exists an element $g\in G$ such that
\[
    |g|=p.
\]
\end{theorem}

\begin{proof}
    Let
    \[
        \Omega
        =
        \{(x_1,x_2,\ldots,x_p)\in G^p : x_1x_2\cdots x_p=1\}.
    \]
    Once \(x_1,\ldots,x_{p-1}\) are chosen, the last entry is uniquely determined by
    \[
        x_p=(x_1x_2\cdots x_{p-1})^{-1}.
    \]
    Hence
    \[
        |\Omega|=|G|^{p-1}.
    \]
    Since \(p\mid |G|\), we have
    \[
        p\mid |\Omega|.
    \]
    
    Let the cyclic group \(C_p=\langle \sigma\rangle\) act on \(\Omega\) by cyclic permutation:
    \[
        \sigma\cdot (x_1,x_2,\ldots,x_p)
        =
        (x_2,x_3,\ldots,x_p,x_1).
    \]
    This is well-defined. Indeed, if \(x_1x_2\cdots x_p=1\), then
    \[
        x_2x_3\cdots x_px_1
        =
        x_1^{-1}(x_1x_2\cdots x_p)x_1
        =
        1.
    \]
    
    The orbits of this action have size either \(1\) or \(p\), because \(C_p\) has prime order.
    Let \(N\) be the number of orbits of size \(1\). Since all remaining orbits have size \(p\), and
    \(p\mid |\Omega|\), it follows that
    \[
        p\mid N.
    \]
    
    Now an element \((x_1,\ldots,x_p)\in\Omega\) is fixed by \(\sigma\) if and only if
    \[
        x_1=x_2=\cdots=x_p.
    \]
    Thus the fixed points are precisely the tuples
    \[
        (x,x,\ldots,x)
    \]
    such that
    \[
        x^p=1.
    \]
    Hence \(N\) is the number of elements \(x\in G\) satisfying \(x^p=1\).
    
    The identity element satisfies \(1^p=1\), so \(N\geq 1\). Since \(p\mid N\), we must have
    \[
        N\geq p.
    \]
    Therefore there exists some \(g\neq 1\) such that
    \[
        g^p=1.
    \]
    Thus \(|g|\) divides \(p\). Since \(p\) is prime and \(g\neq 1\), we conclude that
    \[
        |g|=p.
    \]
    \end{proof}

\begin{proposition}[Orbit decomposition]
Let $G$ act on a set $\Omega$. Then $G$ partitions $\Omega$ into orbits:
\[
    \Omega=\Omega_1\sqcup\Omega_2\sqcup\cdots\sqcup\Omega_m.
\]
\end{proposition}

\begin{proof}
    Define a relation \(\sim\) on \(\Omega\) by
    \[
        \omega\sim \omega'
        \quad\Longleftrightarrow\quad
        \omega'=\omega^g
        \text{ for some } g\in G.
    \]
    We show that \(\sim\) is an equivalence relation.
    
    First, for every \(\omega\in\Omega\), since \(e\in G\) acts trivially, we have
    \[
        \omega^e=\omega.
    \]
    Hence \(\omega\sim\omega\), so \(\sim\) is reflexive.
    
    Next, suppose \(\omega\sim\omega'\). Then there exists \(g\in G\) such that
    \[
        \omega'=\omega^g.
    \]
    Applying \(g^{-1}\), we get
    \[
        \omega
        =
        (\omega')^{g^{-1}}.
    \]
    Thus \(\omega'\sim\omega\), so \(\sim\) is symmetric.
    
    Finally, suppose \(\omega\sim\omega'\) and \(\omega'\sim\omega''\). Then there exist
    \(g,h\in G\) such that
    \[
        \omega'=\omega^g,
        \qquad
        \omega''=(\omega')^h.
    \]
    Therefore
    \[
        \omega''
        =
        (\omega^g)^h
        =
        \omega^{gh}.
    \]
    Since \(gh\in G\), we have \(\omega\sim\omega''\). Hence \(\sim\) is transitive.
    
    Thus \(\sim\) is an equivalence relation on \(\Omega\). Its equivalence classes are exactly
    the orbits of the action, because the equivalence class of \(\omega\) is
    \[
        \{\omega'\in\Omega:\omega'\sim\omega\}
        =
        \{\omega^g:g\in G\}.
    \]
    Equivalence classes of an equivalence relation form a partition of the underlying set.
    Therefore the orbits of \(G\) partition \(\Omega\). If the distinct orbits are denoted by
    \[
        \Omega_1,\Omega_2,\ldots,\Omega_m,
    \]
    then
    \[
        \Omega=\Omega_1\sqcup\Omega_2\sqcup\cdots\sqcup\Omega_m.
    \]
    \end{proof}

\begin{theorem}[Orbit-stabilizer theorem]\label{thm:orbit-stabilizer}
Let $G$ be a finite group acting on a set $\Omega$. If $\Omega_i$ is an orbit and
$\omega_i\in\Omega_i$, then
\[
    |\Omega_i|=|G:G_{\omega_i}|,
\]
where
\[
    G_{\omega_i}
    =
    \{g\in G:\omega_i^g=\omega_i\}
\]
is the stabilizer of $\omega_i$. Equivalently,
\[
    |G|=|\Omega_i|\,|G_{\omega_i}|.
\]
\end{theorem}

\begin{proof}
    Let
    \[
        H=G_{\omega_i}
        =
        \{g\in G:\omega_i^g=\omega_i\}.
    \]
    First we note that \(H\leq G\). Indeed, \(e\in H\), and if \(g,h\in H\), then
    \[
        \omega_i^{gh^{-1}}
        =
        (\omega_i^g)^{h^{-1}}
        =
        \omega_i^{h^{-1}}
        =
        \omega_i,
    \]
    so \(gh^{-1}\in H\). Hence \(H\) is a subgroup of \(G\).
    
    Now consider the set of left cosets \(H\backslash G\). Define a map
    \[
        \varphi:H\backslash G\longrightarrow \Omega_i
    \]
    by
    \[
        \varphi(Hg)=\omega_i^g.
    \]
    
    We first show that \(\varphi\) is well-defined. Suppose
    \[
        Hg=Hh.
    \]
    Then \(h=ag\) for some \(a\in H\). Hence
    \[
        \omega_i^h
        =
        \omega_i^{ag}
        =
        (\omega_i^a)^g
        =
        \omega_i^g,
    \]
    because \(a\in H=G_{\omega_i}\). Thus \(\varphi(Hg)=\varphi(Hh)\).
    
    The map \(\varphi\) is surjective by the definition of the orbit:
    \[
        \Omega_i
        =
        \{\omega_i^g:g\in G\}.
    \]
    
    It remains to show that \(\varphi\) is injective. Suppose
    \[
        \varphi(Hg)=\varphi(Hh).
    \]
    Then
    \[
        \omega_i^g=\omega_i^h.
    \]
    Applying \(h^{-1}\) to both sides gives
    \[
        \omega_i^{gh^{-1}}=\omega_i.
    \]
    Thus
    \[
        gh^{-1}\in H.
    \]
    Hence
    \[
        Hg=Hh.
    \]
    Therefore \(\varphi\) is injective.
    
    Thus \(\varphi\) is a bijection, so
    \[
        |\Omega_i|
        =
        |H\backslash G|
        =
        |G:H|
        =
        |G:G_{\omega_i}|.
    \]
    Since \(G\) is finite, this is equivalent to
    \[
        |G|
        =
        |\Omega_i|\,|G_{\omega_i}|.
    \]
    \end{proof}

\begin{example}[Conjugation action and class equation]
Let $G$ act on itself by conjugation:
\[
    x^g=g^{-1}xg.
\]
For $x\in G$, write
\[
    C(x)=x^G=\{g^{-1}xg:g\in G\}
\]
for the conjugacy class of $x$. Then
\[
    |C(x)|=|G:C_G(x)|,
\]
where
\[
    C_G(x)=\{g\in G:gx=xg\}.
\]
Thus
\[
    |G|=|C(x)|\,|C_G(x)|.
\]
In particular,
\[
    |C(x)|\mid |G|.
\]

If $x_1,\ldots,x_t$ are representatives of the noncentral conjugacy classes of $G$, then
\[
    G
    =
    Z(G)\sqcup C(x_1)\sqcup\cdots\sqcup C(x_t).
\]
Hence the class equation is
\[
    |G|
    =
    |Z(G)|+|C(x_1)|+\cdots+|C(x_t)|.
\]
\end{example}

\begin{theorem}[Center of a finite $p$-group]\label{thm:p-group-center}
Let $G$ be a finite group of order
\[
    |G|=p^e,
\]
where $p$ is prime. Then
\[
    Z(G)\neq 1.
\]
\end{theorem}

\begin{proof}
Use the class equation:
\[
    G
    =
    Z(G)\sqcup C(x_1)\sqcup\cdots\sqcup C(x_t),
\]
where $x_1,\ldots,x_t$ represent the noncentral conjugacy classes. Then
\[
    |G|
    =
    |Z(G)|+|C(x_1)|+\cdots+|C(x_t)|.
\]

For each noncentral element $x_i$, we have
\[
    |C(x_i)|=|G:C_G(x_i)|.
\]
Since $x_i\notin Z(G)$, the class $C(x_i)$ is not a singleton. Therefore
\[
    |C(x_i)|>1.
\]
Also $|C(x_i)|$ divides $|G|=p^e$, so
\[
    |C(x_i)|=p^{e_i}
\]
for some $e_i>0$. Hence each $|C(x_i)|$ is divisible by $p$.

Thus
\[
    |G|
    =
    |Z(G)|+\sum_{i=1}^t |C(x_i)|
\]
implies
\[
    0\equiv |Z(G)| \pmod p.
\]
So $p\mid |Z(G)|$. In particular, $|Z(G)|>1$, and hence
\[
    Z(G)\neq 1.
\]
\end{proof}

\begin{remark}
If $x$ and $y$ are conjugate in $G$, then
\[
    |x|=|y|.
\]
The converse is not true in general. For example, in the cyclic group
\[
    C_m=\langle g\rangle,
\]
the elements $g$ and $g^{-1}$ have the same order, but they are not conjugate unless
$g=g^{-1}$, because $C_m$ is abelian.
\end{remark}

\begin{exercise}
Let $G$ be a finite group with the property that any two elements of the same order are conjugate. Show that
\[
    G\cong S_2
    \qquad\text{or}\qquad
    G\cong S_3.
\]
\end{exercise}

\section{Sylow Theorems}

Throughout this section, let $G$ be a finite group and write
\[
    |G|=p^e m,
\]
where $p$ is prime and
\[
    \gcd(p,m)=1.
\]

\begin{definition}[Sylow $p$-subgroup]
A subgroup $P\leq G$ is called a \textbf{Sylow $p$-subgroup} of $G$ if
\[
    |P|=p^e.
\]
The set of all Sylow $p$-subgroups of $G$ is denoted by
\[
    \operatorname{Syl}_p(G).
\]
\end{definition}

\begin{remark}
If a Sylow $p$-subgroup exists, then Cauchy's theorem follows. Indeed, if
\[
    P\in \operatorname{Syl}_p(G),
\]
then $|P|=p^e$. Choose $1\neq x\in P$. Since $P$ is a $p$-group, $|x|=p^r$ for some $r\geq 1$. Then
\[
    x^{p^{r-1}}
\]
has order $p$.
\end{remark}

\begin{theorem}[First Sylow theorem]\label{thm:first-sylow}
Let $G$ be a finite group. If
\[
    |G|=p^e m,
    \qquad
    \gcd(p,m)=1,
\]
then $G$ has a Sylow $p$-subgroup. That is,
\[
    \operatorname{Syl}_p(G)\neq\varnothing.
\]
\end{theorem}

\begin{proof}
We argue by induction on $|G|$.

Use the class equation:
\[
    |G|
    =
    |Z(G)|+|C(x_1)|+\cdots+|C(x_t)|,
\]
where $x_1,\ldots,x_t$ are representatives of the noncentral conjugacy classes.

\medskip

\noindent\textbf{Case 1: $p\nmid |Z(G)|$.}

Since
\[
    p\mid |G|
\]
but
\[
    p\nmid |Z(G)|,
\]
not all of the noncentral conjugacy class sizes can be divisible by $p$. Hence there exists some $x_i$ such that
\[
    p\nmid |C(x_i)|.
\]
By the orbit-stabilizer theorem,
\[
    |G|=|C(x_i)|\,|C_G(x_i)|.
\]
Since $p^e\mid |G|$ and $p\nmid |C(x_i)|$, we get
\[
    p^e\mid |C_G(x_i)|.
\]
Also $x_i\notin Z(G)$, so
\[
    C_G(x_i)<G.
\]
By induction, $C_G(x_i)$ has a subgroup $P$ of order $p^e$. Then $P$ is a Sylow $p$-subgroup of $G$.

\medskip

\noindent\textbf{Case 2: $p\mid |Z(G)|$.}

First we prove Cauchy's theorem for finite abelian groups.

Let $A$ be a finite abelian group and suppose $p\mid |A|$. We prove that $A$ contains an element of order $p$ by induction on $|A|$.

Choose $1\neq x\in A$. If
\[
    p\mid |x|,
\]
then
\[
    x^{|x|/p}
\]
has order $p$.

Now suppose
\[
    p\nmid |x|.
\]
Then $p$ divides
\[
    |A/\langle x\rangle|.
\]
By induction, $A/\langle x\rangle$ contains an element
\[
    a\langle x\rangle
\]
of order $p$. The order of $a\langle x\rangle$ divides the order of $a$, so
\[
    p\mid |a|.
\]
Hence
\[
    a^{|a|/p}
\]
has order $p$.

Thus every finite abelian group whose order is divisible by $p$ contains an element of order $p$.

Now return to $G$. Since $Z(G)$ is abelian and
\[
    p\mid |Z(G)|,
\]
there exists
\[
    x\in Z(G)
\]
such that
\[
    |x|=p.
\]
Then
\[
    \langle x\rangle\lhd G
\]
and
\[
    |G/\langle x\rangle|
    =
    p^{e-1}m.
\]

If $e=1$, then $\langle x\rangle$ itself is a Sylow $p$-subgroup of $G$.

If $e>1$, then by induction $G/\langle x\rangle$ has a subgroup
\[
    H/\langle x\rangle
\]
of order $p^{e-1}$. Therefore
\[
    |H|=p^e,
\]
so $H$ is a Sylow $p$-subgroup of $G$.
\end{proof}

\begin{theorem}[Second Sylow theorem]\label{thm:second-sylow}
Sylow $p$-subgroups are conjugate. More generally, if $H\leq G$ is a $p$-subgroup and
\[
    P\in\operatorname{Syl}_p(G),
\]
then there exists $x\in G$ such that
\[
    H\leq P^x.
\]
In particular, if $H$ is also a Sylow $p$-subgroup, then
\[
    H=P^x.
\]
\end{theorem}

\begin{proof}
Let
\[
    \Delta=\{P^x:x\in G\}
\]
be the conjugacy class of $P$ under conjugation. Then $G$ acts transitively on $\Delta$ by
\[
    (P^x)^g=P^{xg}.
\]
By the orbit-stabilizer theorem,
\[
    |\Delta|
    =
    |G:N_G(P)|.
\]
Since
\[
    P\leq N_G(P),
\]
we have
\[
    |\Delta|=|G:N_G(P)|\mid |G:P|=m.
\]
Hence
\[
    p\nmid |\Delta|.
\]

Now let $H\leq G$ be a $p$-subgroup, say
\[
    |H|=p^d.
\]
Restrict the action of $G$ on $\Delta$ to an action of $H$. Then $\Delta$ is a disjoint union of $H$-orbits:
\[
    \Delta=\Delta_1\sqcup\cdots\sqcup\Delta_r.
\]
Each orbit size divides $|H|$, so each $|\Delta_i|$ is a power of $p$. Since
\[
    p\nmid |\Delta|,
\]
at least one orbit has size $1$.

Thus there exists $x\in G$ such that
\[
    (P^x)^h=P^x
    \qquad
    \text{for every }h\in H.
\]
Equivalently,
\[
    H\leq N_G(P^x).
\]

But $P^x$ is a normal Sylow $p$-subgroup of $N_G(P^x)$. Therefore every $p$-subgroup of $N_G(P^x)$ is contained in $P^x$. Hence
\[
    H\leq P^x.
\]

If $H$ is itself a Sylow $p$-subgroup, then
\[
    |H|=|P^x|=p^e,
\]
so
\[
    H=P^x.
\]
Therefore all Sylow $p$-subgroups are conjugate.
\end{proof}

\begin{theorem}[Third Sylow theorem]\label{thm:third-sylow}
Let
\[
    |G|=p^e m,
    \qquad
    \gcd(p,m)=1,
\]
and let
\[
    n_p=|\operatorname{Syl}_p(G)|
\]
be the number of Sylow $p$-subgroups of $G$. Then
\[
    n_p\mid m
\]
and
\[
    n_p\equiv 1\pmod p.
\]
\end{theorem}

\begin{proof}
Let
\[
    P\in\operatorname{Syl}_p(G).
\]
By the second Sylow theorem,
\[
    \operatorname{Syl}_p(G)
    =
    \{P^x:x\in G\}.
\]
Therefore
\[
    n_p
    =
    |G:N_G(P)|.
\]
Since
\[
    P\leq N_G(P),
\]
we get
\[
    n_p=|G:N_G(P)|\mid |G:P|=m.
\]
Hence
\[
    n_p\mid m.
\]

Now let $P$ act on $\operatorname{Syl}_p(G)$ by conjugation. The orbit sizes are powers of $p$. We determine the fixed points.

Suppose
\[
    Q\in \operatorname{Syl}_p(G)
\]
is fixed by $P$. Then
\[
    Q^h=Q
    \qquad
    \text{for every }h\in P,
\]
so
\[
    P\leq N_G(Q).
\]
But $Q$ is a normal Sylow $p$-subgroup of $N_G(Q)$. Hence every $p$-subgroup of $N_G(Q)$ is contained in $Q$. Thus
\[
    P\leq Q.
\]
Since $P$ and $Q$ have the same order, we get
\[
    P=Q.
\]
Therefore the only fixed point of this action is $P$ itself.

All other $P$-orbits have sizes divisible by $p$. Hence
\[
    n_p
    =
    1+\text{(a sum of multiples of }p),
\]
so
\[
    n_p\equiv 1\pmod p.
\]
\end{proof}

